% !TeX program = pdflatex
% Podręcznik podejmowania decyzji --- Basilisk Systems P.S.A.
% Streszcza umowę spółki (psa.tex), umowę akcjonariuszy (basilisk_sha.tex) i dokumenty wewnętrzne.
\documentclass[11pt,a4paper]{report}
\usepackage[T1]{fontenc}
\usepackage[utf8]{inputenc}
\usepackage[english,polish]{babel}
\usepackage{basilisk_i18n}
\usepackage{lmodern}
\usepackage{microtype}
\usepackage{setspace}
\usepackage{parskip}
\usepackage{needspace}
\usepackage{tabularx}
\usepackage{float}
\usetikzlibrary{positioning,arrows.meta,shapes.geometric,calc,fit}

\BusinessPlanCompany{Basilisk Systems}
\BusinessPlanWatermarkText{Basilisk Systems}
\BusinessPlanWatermarkColor{black!6}
\BusinessPlanWatermarkFontSize{38}{46}
\BusinessPlanWatermarkAngle{38}
\BusinessPlanAuthorsText{Basilisk Systems --- Założyciele}
\BusinessPlanConfidentialText{Dokument wewnętrzny}
\BusinessPlanFooterText{Basilisk Systems --- podręcznik podejmowania decyzji}
\BusinessPlanVersionText{Wersja 1.0 --- projekt}

\setstretch{1.1}
\setlist[itemize]{topsep=0.2em,itemsep=0.12em,parsep=0.05em,leftmargin=1.8em}
\setlist[enumerate]{topsep=0.2em,itemsep=0.12em,parsep=0.05em,leftmargin=2.0em}
\setlength{\parindent}{0pt}
\setlength{\parskip}{0.5em}
\setlength{\emergencystretch}{4em}
\renewcommand{\arraystretch}{1.2}
\renewcommand{\thesection}{\arabic{section}}
\captionsetup{font=small,labelfont=bf}

\hypersetup{
  pdftitle={Basilisk Systems P.S.A. --- podręcznik podejmowania decyzji},
  pdfauthor={Basilisk Systems}
}

% ---------------------------------------------------------------- style diagramów
\definecolor{cWZ}{HTML}{1F4E79}
\definecolor{cRD}{HTML}{2E7D32}
\definecolor{cDW}{HTML}{8D6E00}
\definecolor{cX}{HTML}{8E2430}
\tikzset{
  >={Stealth[length=2.2mm]},
  every node/.style={font=\small, align=center},
  box/.style={draw, rounded corners=2pt, minimum height=8mm, inner sep=4pt, text width=#1},
  box/.default=3.6cm,
  wz/.style={box=#1, fill=cWZ!12, draw=cWZ},
  wz/.default=3.6cm,
  rd/.style={box=#1, fill=cRD!12, draw=cRD},
  rd/.default=3.6cm,
  dw/.style={box=#1, fill=cDW!14, draw=cDW},
  dw/.default=3.6cm,
  xx/.style={box=#1, fill=cX!10, draw=cX},
  xx/.default=3.6cm,
  q/.style={draw, diamond, aspect=2.4, inner sep=1.5pt, fill=black!4, text width=#1},
  q/.default=3.4cm,
  lbl/.style={font=\scriptsize, fill=white, inner sep=1.5pt},
  line/.style={->, thick},
}
\newtcolorbox{Uwaga}{colback=cX!5, colframe=cX!70, boxrule=0.6pt, arc=2pt, left=6pt, right=6pt, top=4pt, bottom=4pt}
\newtcolorbox{Zasada}{colback=cWZ!5, colframe=cWZ!70, boxrule=0.6pt, arc=2pt, left=6pt, right=6pt, top=4pt, bottom=4pt}

\begin{document}

\BusinessPlanTitlePage
  {Basilisk Systems}
  {Wersja 1.0 --- projekt}
  {Podręcznik podejmowania decyzji}
  {Kto decyduje, jaką większością i w jakim trybie}
  {Walne Zgromadzenie, Rada Dyrektorów, dyrektor wykonawczy --- progi, terminy, procedury}
  {1.0}

\tableofcontents
\newpage

\section{Jak korzystać z podręcznika}

Podręcznik streszcza zasady podejmowania decyzji wynikające z umowy spółki (dalej „US”, wersja 3.3-S), umowy akcjonariuszy (dalej „UA”, wersja 1.7) oraz dokumentów wewnętrznych: polityki bezpieczeństwa informacji, Cennika i wzoru NDA. Odwołania w nawiasach wskazują źródło zasady.

\begin{Uwaga}
Podręcznik nie zmienia żadnej umowy. W razie rozbieżności rozstrzyga treść umowy spółki, a następnie umowy akcjonariuszy. Limity wewnętrzne wiążą dyrektorów wobec Spółki, ale nie wpływają na ważność umów z osobami trzecimi (US §~11 ust.~4): dyrektor, który je przekroczy, odpowiada wobec Spółki.
\end{Uwaga}

Kolory na diagramach: \textcolor{cWZ}{\textbf{niebieski}} --- akcjonariusze i Walne Zgromadzenie; \textcolor{cRD}{\textbf{zielony}} --- Rada Dyrektorów; \textcolor{cDW}{\textbf{złoty}} --- dyrektor wykonawczy; \textcolor{cX}{\textbf{czerwony}} --- sytuacje szczególne (spór, sąd, konflikt interesów).

% ================================================================
\section{Organy Spółki}

Spółka działa w systemie monistycznym: ma tylko Walne Zgromadzenie i Radę Dyrektorów; nie ma zarządu ani rady nadzorczej (US §~10 ust.~1).

\begin{figure}[H]
\centering
\begin{tikzpicture}[node distance=9mm and 6mm]
  \node[wz=5.2cm] (akc) {\textbf{Akcjonariusze}\\1 akcja = 1 głos};
  \node[wz=5.2cm, below=of akc] (wz) {\textbf{Walne Zgromadzenie}\\albo uchwała poza WZ (e-mail)};
  \node[rd=5.2cm, below=14mm of wz] (rd) {\textbf{Rada Dyrektorów}\\3--7 dyrektorów, kadencja 3 lata\\Przewodniczący: głos rozstrzygający};
  \node[dw=3.8cm, below=14mm of rd, xshift=-2.6cm] (dw) {\textbf{Dyrektor wykonawczy}\\(co najmniej jeden)\\bieżąca działalność};
  \node[rd=3.8cm, below=14mm of rd, xshift=2.6cm] (dn) {\textbf{Dyrektor niewykonawczy}\\(co najmniej jeden)\\nadzór, umowy z dyr.~wyk.};
  \node[box=2.6cm, right=22mm of rd] (pr) {\textbf{Prokurent}\\(opcjonalnie)};
  \draw[line] (akc) -- (wz);
  \draw[line] (wz) -- node[lbl, right, text width=4.2cm, align=left]{powołuje i odwołuje\\(bezwzględna większość oddanych)} (rd);
  \draw[line] (rd.south) ++(-1.2,0) -- node[lbl, left]{deleguje} (dw.north);
  \draw[line] (rd.south) ++(1.2,0) -- (dn.north);
  \draw[line] (rd) -- node[lbl, above, text width=1.9cm]{zgoda wszystkich dyrektorów} (pr);
  \node[box=10cm, below=8mm of dw.south east, anchor=north, xshift=0.6cm, fill=black!3] (rep) {\textbf{Reprezentacja wobec osób trzecich:} dwóch dyrektorów łącznie albo dyrektor z prokurentem (US §~11 ust.~1)};
\end{tikzpicture}
\caption{Struktura organów (US §~10--12).}
\end{figure}

\begin{itemize}
  \item \textbf{Rada Dyrektorów} prowadzi sprawy Spółki, reprezentuje ją i nadzoruje (US §~10 ust.~2). Uchwały zapadają bezwzględną większością głosów przy udziale co najmniej połowy dyrektorów; przy równości głosów decyduje Przewodniczący, chyba że sprawa dotyczy jego interesu (US §~10 ust.~5).
  \item \textbf{Dyrektor wykonawczy} prowadzi bieżącą działalność na podstawie delegacji Rady. Delegacja nie obejmuje spraw wymagających uchwały Rady (US §~10 ust.~4 i 7).
  \item \textbf{Dyrektor niewykonawczy} --- musi być co najmniej jeden. Reprezentuje Spółkę w umowach i sporach z dyrektorem wykonawczym na podstawie uchwały dyrektorów niewykonawczych (US §~11 ust.~3).
  \item \textbf{Regulamin Rady Dyrektorów} (przyjmuje Rada) określa szczegółowy podział kompetencji, limity wydatków i zasady raportowania (US §~10 ust.~8).
\end{itemize}

% ================================================================
\section{Kto decyduje? --- drzewo decyzyjne}

Przy każdej nowej sprawie przejdź drzewo od góry. Pierwsza odpowiedź „tak” wskazuje organ i większość.

\begin{figure}[H]
\centering
\begin{tikzpicture}[node distance=7mm and 9mm]
  \node[box=3.2cm, fill=black!8] (start) {\textbf{Nowa sprawa}};
  \node[q, below=of start] (q1) {Pozbawienie prawa poboru?};
  \node[wz=4.4cm, right=of q1] (a1) {\textbf{WZ --- 4/5 głosów}\\US §~5 ust.~6};
  \node[q, below=of q1] (q2) {Sprawa z listy\\„75\%” (tabela 1)?};
  \node[wz=4.4cm, right=of q2] (a2) {\textbf{WZ --- 75\% wszystkich głosów w Spółce}\\US §~13 ust.~1};
  \node[q, below=of q2] (q3) {Dyrektorzy, sprawozdanie, absolutorium, zysk, wynagrodzenie dyrektorów?};
  \node[wz=4.4cm, right=of q3] (a3) {\textbf{WZ --- bezwzględna większość oddanych}\\US §~13 ust.~3};
  \node[q, below=of q3] (q4) {Sprawa z listy Rady (tabela 2)?};
  \node[rd=4.4cm, right=of q4] (a4) {\textbf{Uchwała Rady Dyrektorów}\\US §~10 ust.~5 i 7};
  \node[dw=4.4cm, below=of q4] (a5) {\textbf{Dyrektor wykonawczy}\\w granicach budżetu i regulaminu Rady};
  \draw[line] (start) -- (q1);
  \foreach \i/\j in {1/2,2/3,3/4} \draw[line] (q\i) -- node[lbl, left]{nie} (q\j);
  \draw[line] (q4) -- node[lbl, left]{nie} (a5);
  \foreach \i in {1,2,3,4} \draw[line] (q\i) -- node[lbl, above]{tak} (a\i);
\end{tikzpicture}
\caption{Drzewo decyzyjne: który organ i jaka większość.}
\end{figure}

\begin{Zasada}
\textbf{„75\% wszystkich głosów w Spółce”} liczy się od wszystkich istniejących głosów, a nie od obecnych. Nieobecny albo wstrzymujący się akcjonariusz działa więc jak głos „przeciw”. \textbf{„Bezwzględna większość oddanych”} oznacza więcej głosów „za” niż łącznie „przeciw” i „wstrzymujących się” spośród oddanych (US §~12 ust.~7).
\end{Zasada}

\needspace{12\baselineskip}
\subsection*{Tabela 1. Sprawy zastrzeżone dla Walnego Zgromadzenia --- 75\% wszystkich głosów (US §~13 ust.~1)}
\begin{tabularx}{\linewidth}{@{}p{0.6cm} X@{}}
\toprule
a) & zmiana umowy spółki (na Walnym Zgromadzeniu, z protokołem notarialnym, i wpis do rejestru) \\
b) & emisja akcji, warrantów, opcji i innych praw do akcji, w tym każda emisja serii F i P \\
c) & nadanie albo cofnięcie statusu Założyciela Funkcjonalnego Serii~F \\
d) & połączenie, podział, przekształcenie, rozwiązanie, likwidacja \\
e) & zbycie albo dzierżawa przedsiębiorstwa lub jego zorganizowanej części; zbycie aktywów ponad 25\% wartości aktywów \\
f) & przeniesienie albo wyłączna, nieodwołalna lub bezterminowa licencja na kluczową własność intelektualną poza zwykłym tokiem działalności \\
g) & nabycie akcji własnych, umorzenie, wypłata z kapitału akcyjnego \\
h) & zobowiązanie, zabezpieczenie albo transakcja poza budżetem powyżej 500\,000 zł \\
i) & zgoda na zbycie akcji w okresie lock-up (12 miesięcy) i zwolnienie akcji z obowiązku zbycia ponad zwolnienie automatyczne \\
j) & kredyt, pożyczka, poręczenie lub podobna umowa z dyrektorem, prokurentem, likwidatorem albo na ich rzecz \\
\bottomrule
\end{tabularx}

\needspace{12\baselineskip}
\subsection*{Tabela 2. Sprawy wymagające uchwały Rady Dyrektorów (US §~10 ust.~7)}
\begin{tabularx}{\linewidth}{@{}p{0.6cm} X@{}}
\toprule
1 & przyjęcie i zmiana rocznego budżetu i planu finansowego \\
2 & umowa albo wydatek poza budżetem powyżej 250\,000 zł \\
3 & finansowanie, zabezpieczenie albo rozporządzenie aktywami powyżej 100\,000 zł \\
4 & transakcja z akcjonariuszem, dyrektorem, osobą bliską albo podmiotem powiązanym \\
5 & utworzenie spółki zależnej, działalność w nowej jurysdykcji \\
6 & licencja wyłączna albo przeniesienie istotnej własności intelektualnej (jeżeli nie wymaga WZ --- tabela~1 lit.~f) \\
7 & przyjęcie polityk: bezpieczeństwa informacji, kontroli eksportu, zgodności; Cennika \\
8 & wykonanie uchwał emisyjnych; stwierdzenie wniesienia wkładów \\
9 & zgoda albo odmowa zgody na zbycie lub obciążenie akcji (rozdział~6) \\
10 & stwierdzenie ustania zaangażowania i wskazanie nabywcy akcji niezwolnionych (rozdział~7) \\
11 & pożyczka od Założyciela (UA §~10 ust.~2) \\
\bottomrule
\end{tabularx}

% ================================================================
\section{Progi kwotowe}

Ta sama czynność może wymagać różnych decyzji zależnie od kwoty. Kwoty dotyczą jednej czynności albo kilku czynności powiązanych.

\begin{figure}[H]
\centering
\begin{tikzpicture}[x=0.0165cm, y=1cm]
  % oś
  \draw[->] (0,0) -- (640,0) node[right, font=\scriptsize]{tys. zł};
  \foreach \v in {0,100,250,500} {\draw (\v,0.08) -- (\v,-0.08) node[below, font=\scriptsize]{\v};}
  % pas 1: wydatek poza budżetem
  \node[anchor=east, font=\scriptsize, text width=3.4cm, align=right] at (-8,1.2) {umowa / wydatek\\poza budżetem};
  \fill[cDW!25] (0,0.85) rectangle (250,1.55);  \node[font=\scriptsize] at (125,1.2) {dyrektor wykonawczy};
  \fill[cRD!25] (250,0.85) rectangle (500,1.55); \node[font=\scriptsize] at (375,1.2) {Rada Dyrektorów};
  \fill[cWZ!25] (500,0.85) rectangle (630,1.55); \node[font=\scriptsize] at (565,1.2) {WZ 75\%};
  % pas 2: finansowanie / zabezpieczenie / aktywa
  \node[anchor=east, font=\scriptsize, text width=3.4cm, align=right] at (-8,2.3) {finansowanie, zabezpie-\\czenie, zbycie aktywów};
  \fill[cDW!25] (0,1.95) rectangle (100,2.65);  \node[font=\scriptsize] at (50,2.3) {dyr.~wyk.};
  \fill[cRD!25] (100,1.95) rectangle (500,2.65); \node[font=\scriptsize] at (300,2.3) {Rada Dyrektorów};
  \fill[cWZ!25] (500,1.95) rectangle (630,2.65); \node[font=\scriptsize] at (565,2.3) {WZ 75\%};
  % pas 3: w budżecie
  \node[anchor=east, font=\scriptsize, text width=3.4cm, align=right] at (-8,3.4) {wydatek w ramach\\przyjętego budżetu};
  \fill[cDW!25] (0,3.05) rectangle (630,3.75); \node[font=\scriptsize] at (315,3.4) {dyrektor wykonawczy (w granicach regulaminu Rady)};
  \foreach \v in {100,250,500} \draw[densely dotted, black!50] (\v,0) -- (\v,3.85);
\end{tikzpicture}
\caption{Kto zatwierdza czynność w zależności od kwoty (US §~10 ust.~7, §~13 ust.~1 lit.~h).}
\end{figure}

Niezależnie od kwoty: transakcja z akcjonariuszem, dyrektorem albo podmiotem powiązanym wymaga uchwały Rady (bez udziału dyrektora mającego konflikt interesów), a umowa kredytu lub pożyczki z dyrektorem --- także zgody WZ 75\% (tabela~1 lit.~j). Zbycie aktywów ponad 25\% ich wartości zawsze wymaga WZ 75\% (lit.~e).

% ================================================================
\section{Jak podjąć uchwałę}

\subsection*{Uchwała akcjonariuszy poza Walnym Zgromadzeniem (e-mail)}

Domyślny tryb dla większości spraw. Wyjątek: uchwały wymagające protokołu notarialnego (w szczególności zmiana umowy spółki) podejmuje się na Walnym Zgromadzeniu (US §~12 ust.~1).

\begin{figure}[H]
\centering
\begin{tikzpicture}[node distance=6mm and 5mm]
  \node[wz=2.9cm] (s1) {\textbf{1. Projekt}\\Rada, Przewodniczący albo akcjonariusz z $\ge$5\% głosów rozsyła projekt z uzasadnieniem};
  \node[wz=2.9cm, right=of s1] (s2) {\textbf{2. Głosowanie}\\e-mailem z adresu z rejestru; termin: 3 dni robocze -- 21 dni};
  \node[wz=2.9cm, right=of s2] (s3) {\textbf{3. Podjęcie}\\z upływem terminu albo wcześniej, gdy jest większość wszystkich głosów};
  \node[wz=2.9cm, right=of s3] (s4) {\textbf{4. Wynik}\\Przewodniczący stwierdza wynik i zawiadamia e-mailem; wpis do księgi protokołów};
  \draw[line] (s1) -- (s2); \draw[line] (s2) -- (s3); \draw[line] (s3) -- (s4);
\end{tikzpicture}
\caption{Uchwała poza Walnym Zgromadzeniem (US §~12 ust.~2--3 i 8).}
\end{figure}

\begin{itemize}
  \item Krótszy termin niż 3 dni robocze jest możliwy za zgodą wszystkich akcjonariuszy. Głos oddany po terminie się nie liczy.
  \item Głosowanie \emph{na piśmie} (papier) wymaga zgody wszystkich akcjonariuszy na ten tryb; głosowanie e-mailem takiej zgody nie wymaga.
  \item Walne Zgromadzenie bez formalnego zwołania jest możliwe, jeżeli reprezentowane są wszystkie akcje i nikt nie zgłasza sprzeciwu (US §~12 ust.~4). Udział zdalny reguluje regulamin WZ.
\end{itemize}

\subsection*{Uchwała Rady Dyrektorów}

\begin{enumerate}
  \item Przewodniczący (albo dyrektor) rozsyła projekt; głosowanie może odbyć się na posiedzeniu, na piśmie albo e-mailem, jeżeli da się zidentyfikować uczestników i utrwalić głosy (US §~10 ust.~6).
  \item \textbf{Konflikt interesów:} dyrektor, którego interes jest sprzeczny z interesem Spółki, ujawnia to i nie głosuje; nie liczy się go do kworum.
  \item Kworum: co najmniej połowa powołanych dyrektorów. Większość: bezwzględna. Przy równości --- głos Przewodniczącego.
  \item Uchwałę zapisuje się i przechowuje w dokumentacji Spółki.
\end{enumerate}

% ================================================================
\section{Zgoda na zbycie akcji}

Każde zbycie albo obciążenie akcji wymaga zgody Spółki, którą wydaje Rada Dyrektorów (US §~8). Spółka \textbf{musi} wyrazić zgodę, jeżeli nie zachodzi żadna z pięciu przyczyn odmowy.

\begin{figure}[H]
\centering
\begin{tikzpicture}[node distance=7mm and 8mm]
  \node[box=4.2cm, fill=black!8] (z) {\textbf{Zgłoszenie do Spółki}\\nabywca, cena, istotne warunki};
  \node[q=3.3cm, below=of z] (lock) {Lock-up 12 mies.\\od wpisu / Dnia Przyznania F?};
  \node[wz=3.4cm, right=of lock] (lockwz) {zgoda WZ 75\%\\(albo czekać)};
  \node[rd=4.2cm, below=of lock] (rd) {\textbf{Rada Dyrektorów}\\14 dni na zgodę albo odmowę\\(forma dokumentowa, e-mail)};
  \node[rd=3.4cm, below=10mm of rd] (no) {\textbf{Odmowa}\\tylko z przyczyn a)--e), z ich wskazaniem};
  \node[rd=2.8cm, left=6mm of no] (ok) {\textbf{Zgoda}\\$\to$ wpis nabywcy do rejestru};
  \node[xx=3.2cm, right=6mm of no] (cisza) {\textbf{Milczenie}\\$\ne$ zgoda};
  \node[rd=3.4cm, below=of no] (e) {Przyczyna e) (konkurent): Spółka w 30 dni wskazuje innego nabywcę na tych samych warunkach --- inaczej musi udzielić zgody};
  \node[xx=3.2cm, below=of cisza] (w7) {wezwanie Rady\\+ 7 dni};
  \node[wz=3.2cm, below=of w7] (ua) {uchwała akcjonariuszy poza WZ --- 14 dni (tylko z przyczyn a--e)};
  \node[xx=3.2cm, below=of ua] (sad) {sąd: żądanie złożenia oświadczenia o zgodzie (art.~64 KC)};
  \draw[line] (z) -- (lock);
  \draw[line] (lock) -- node[lbl, above]{tak} (lockwz);
  \draw[line] (lock) -- node[lbl, left]{nie} (rd);
  \draw[line] (rd) -- (ok); \draw[line] (rd) -- (no); \draw[line] (rd) -- (cisza);
  \draw[line] (no) -- (e);
  \draw[line] (cisza) -- (w7); \draw[line] (w7) -- (ua); \draw[line] (ua) -- node[lbl, right]{brak zgody\\i brak odmowy} (sad);
\end{tikzpicture}
\caption{Procedura zgody na zbycie akcji (US §~8 ust.~1--4).}
\end{figure}

\textbf{Przyczyny odmowy (zamknięty katalog, US §~8 ust.~2):} a)~sankcje albo sprzeczność z przepisami o kontroli inwestycji, eksportu lub bezpieczeństwie państwa; b)~brak wymaganej zgody organu (można dać zgodę warunkową); c)~lock-up bez zgody WZ; d)~akcje objęte obowiązkiem zbycia, a nabywca nie przejął obowiązków (pełnomocnictwo i oferta); e)~nabywca jest konkurentem (z wyjątkiem inwestycji portfelowych do 5\%).

Zgody nie wymaga nabycie akcji w wykonaniu obowiązku zbycia ani zbycie w egzekucji (US §~8 ust.~5).

% ================================================================
\section{Vesting i ustanie zaangażowania}

\subsection*{Harmonogram zwalniania akcji (US §~9 ust.~2)}

\begin{figure}[H]
\centering
\begin{tikzpicture}[x=0.245cm, y=0.042cm]
  \draw[->] (0,0) -- (51,0) node[right, font=\scriptsize]{miesiące};
  \draw[->] (0,0) -- (0,108) node[above, font=\scriptsize]{\% akcji zwolnionych};
  \foreach \m in {0,6,12,24,36,48} \draw (\m,1.5) -- (\m,-1.5) node[below, font=\scriptsize]{\m};
  \foreach \p in {25,50,75,100} {\draw (0.3,\p) -- (-0.3,\p) node[left, font=\scriptsize]{\p}; \draw[black!12] (0,\p) -- (50,\p);}
  % Założyciel Pierwotny
  \draw[very thick, cWZ] (0,25) -- (6,25) -- (6,37.5) -- (36,100) -- (49,100);
  \fill[cWZ] (0,25) circle (0.25 and 1.6);
  \node[font=\scriptsize, cWZ, anchor=west, text width=4.2cm, align=left] at (1,90) {\textbf{Założyciel Pierwotny:} 25\% od razu, cliff 6 mies., potem co miesiąc do 36 mies.};
  % Seria F
  \draw[very thick, cRD, dashed] (0,0) -- (12,0) -- (12,25) -- (48,100) -- (49,100);
  \node[font=\scriptsize, cRD, anchor=west, text width=4.2cm, align=left] at (26,40) {\textbf{Seria F:} cliff 12 mies. (25\%), potem co miesiąc do 48 mies.};
\end{tikzpicture}
\caption{Zwalnianie akcji z obowiązku zbycia. Do chwili zwolnienia akcje dają pełne prawa (US §~9 ust.~3).}
\end{figure}

\subsection*{Co się dzieje przy odejściu (US §~9 ust.~4--5, UA §~4--5)}

\begin{figure}[H]
\centering
\begin{tikzpicture}[node distance=7mm and 6mm]
  \node[box=4.6cm, fill=black!8] (ev) {\textbf{Ustanie zaangażowania}\\(rozwiązanie umowy, rezygnacja, śmierć, niezdolność)};
  \node[rd=4.6cm, below=of ev] (st) {\textbf{Rada stwierdza ustanie} w 30 dni (bez udziału zainteresowanego): dzień i liczba akcji niezwolnionych};
  \node[q=3.1cm, below=of st] (typ) {Rodzaj odejścia?};
  \node[wz=3.6cm, below=10mm of typ] (u) {\textbf{Rozwiązanie bez przyczyny / odejście uzgodnione}\\zwalnia się dodatkowo 12 mies.};
  \node[wz=3.4cm, left=6mm of u] (d) {\textbf{Śmierć / trwała niezdolność}\\zwalniają się \textbf{wszystkie} akcje};
  \node[xx=3.4cm, right=6mm of u] (z) {\textbf{Odejście Zwykłe}\\bez dodatkowego zwolnienia};
  \node[rd=6.2cm, below=12mm of u] (n) {Akcje niezwolnione $\to$ zbycie \textbf{po cenie emisyjnej}: Rada wskazuje nabywcę w 3 mies.; potem Założyciele Pierwotni w 30 dni; po 6 mies. obowiązek wygasa};
  \node[wz=6.2cm, below=of n] (wz75) {WZ może (75\%) zwolnić także pozostałe akcje};
  \draw[line] (ev) -- (st); \draw[line] (st) -- (typ);
  \draw[line] (typ) -| (d); \draw[line] (typ) -- (u); \draw[line] (typ) -| (z);
  \draw[line] (u) -- (n); \draw[line] (z) |- (n);
  \draw[line] (n) -- (wz75);
\end{tikzpicture}
\caption{Ustanie zaangażowania i los akcji niezwolnionych.}
\end{figure}

\begin{itemize}
  \item Zwolnienie przy śmierci, niezdolności i Odejściu Uzasadnionym następuje \textbf{automatycznie} z mocy umowy spółki --- nie wymaga głosowania (UA §~4 ust.~4).
  \item Umowę zbycia w imieniu odchodzącego może zawrzeć pełnomocnik (Spółka albo dyrektor niewykonawczy) na podstawie nieodwołalnego pełnomocnictwa; nabywca płaci w 14 dni, wpis do rejestru dopiero po zapłacie (UA §~5).
  \item Po Odejściu Zwykłym obowiązuje zakaz konkurencji przez 12 miesięcy; zakaz aktywnego nakłaniania pracowników i klientów obowiązuje 12 miesięcy przy każdym odejściu (UA §~8).
\end{itemize}

% ================================================================
\section{Umowy i spory z dyrektorami}

\begin{figure}[H]
\centering
\begin{tikzpicture}[node distance=7mm and 8mm]
  \node[box=4.4cm, fill=black!8] (u) {\textbf{Umowa albo spór Spółki z dyrektorem}};
  \node[q=3.0cm, below=of u] (q) {Dyrektor wykonawczy?};
  \node[rd=4.6cm, below left=9mm and -10mm of q] (dn) {Dyrektor niewykonawczy na podstawie uchwały dyrektorów niewykonawczych\\\textbf{albo} pełnomocnik z uchwały akcjonariuszy};
  \node[wz=4.6cm, below right=9mm and -10mm of q] (pe) {\textbf{Pełnomocnik} powołany uchwałą akcjonariuszy (może być poza WZ; nie może być stroną umowy)};
  \node[xx=9.8cm, below=30mm of q] (poz) {\textbf{Pożyczka, kredyt, poręczenie z dyrektorem:} uchwała Rady + zgoda WZ 75\% + pełnomocnik z uchwały akcjonariuszy (US §~13 ust.~1 lit.~j, UA §~10 ust.~2)};
  \draw[line] (u) -- (q);
  \draw[line] (q) -- node[lbl, left]{tak} (dn);
  \draw[line] (q) -- node[lbl, right]{nie} (pe);
\end{tikzpicture}
\caption{Reprezentacja Spółki w sprawach z dyrektorem (US §~11 ust.~3).}
\end{figure}

Dotyczy to w szczególności umów przeniesienia praw IP (Umowa nr 1 i nr 2) z Założycielami będącymi dyrektorami, umów wykonawczych do obowiązku zbycia i potwierdzeń kwartalnych. Jedna uchwała akcjonariuszy może powołać pełnomocnika do wszystkich umów danego rodzaju.

% ================================================================
\section{Emisje akcji}

\begin{figure}[H]
\centering
\begin{tikzpicture}[node distance=7mm and 6mm]
  \node[q=3.4cm] (q) {Seria A, F albo P,\\w limitach?};
  \node[wz=4.0cm, right=10mm of q] (zm) {\textbf{Najpierw zmiana umowy} (WZ 75\%, notariusz, KRS) --- nowa seria albo wyższy limit};
  \node[wz=4.0cm, below=of q] (em) {\textbf{Uchwała emisyjna WZ 75\%}\\maksymalna liczba akcji, cena minimalna, termin};
  \node[wz=3.4cm, left=10mm of em] (pp) {Wyłączenie prawa poboru?\\$\to$ \textbf{4/5 głosów}};
  \node[rd=4.0cm, below=of em] (rd) {\textbf{Rada wykonuje:} oferty, umowy objęcia, przyjęcie wkładów, rejestr, zgłoszenia};
  \draw[line] (q) -- node[lbl, above]{nie} (zm);
  \draw[line] (q) -- node[lbl, left]{tak} (em);
  \draw[line] (zm) |- (em);
  \draw[line, dashed] (em) -- (pp);
  \draw[line] (em) -- (rd);
\end{tikzpicture}
\caption{Emisja akcji (US §~5--7).}
\end{figure}

\begin{itemize}
  \item \textbf{Limity:} łącznie 2\,000\,000 akcji; seria F do 250\,000 (i do 20\% akcji); seria P do 176\,471; termin emisji do 31 grudnia 2036 r. (US §~5 ust.~2--3).
  \item \textbf{Seria F --- dwie uchwały WZ 75\%:} najpierw nadanie statusu Założyciela Funkcjonalnego (rola, zaangażowanie, maksymalna liczba akcji, Dzień Przyznania), potem emisja (US §~6).
  \item \textbf{Seria P:} regulamin programu zatwierdza WZ; Rada prowadzi program i rekomenduje uczestników, ale nie emituje akcji samodzielnie (US §~7).
  \item Cena emisyjna nie niższa niż 0,01 zł; Rada nie może jej obniżyć poniżej minimum z uchwały (US §~5 ust.~5 i 7).
\end{itemize}

% ================================================================
\section{Dokumenty wewnętrzne --- kto przyjmuje}

\begin{tabularx}{\linewidth}{@{}>{\raggedright\arraybackslash}p{4.0cm} >{\raggedright\arraybackslash}p{3.3cm} X@{}}
\toprule
\textbf{Dokument} & \textbf{Przyjmuje} & \textbf{Uwagi} \\
\midrule
Roczny budżet i plan finansowy & Rada Dyrektorów & wyznacza granice działania dyrektora wykonawczego \\
Regulamin Rady Dyrektorów & Rada Dyrektorów & podział kompetencji, limity, raportowanie, konflikt interesów \\
Regulamin Walnego Zgromadzenia & akcjonariusze & udział zdalny; pierwszy przyjmują Założyciele \\
Polityka bezpieczeństwa informacji & Rada Dyrektorów & wyznacza Odpowiedzialnego; każdy Użytkownik potwierdza zapoznanie się \\
Cennik wynagrodzeń za rozwiązania projektowe & Rada Dyrektorów & waloryzacja GUS co roku 1 marca; zmiana na niekorzyść nie działa wstecz \\
Regulamin programu motywacyjnego (seria P) & Walne Zgromadzenie & warunki, ograniczenia zbywania, skutki ustania współpracy \\
Wzór NDA & Rada Dyrektorów & zawiera się \textbf{przed} ujawnieniem informacji (UA §~7 ust.~2 lit.~b); podpisuje dwóch dyrektorów albo dyrektor z prokurentem \\
Lista osób z dostępem do informacji Kontrolowanych & Rada Dyrektorów & polityka §~3; przegląd dostępów co kwartał \\
\bottomrule
\end{tabularx}

% ================================================================
\section{Impas i spory}

\begin{figure}[H]
\centering
\begin{tikzpicture}[node distance=6mm and 6mm]
  \node[box=3.0cm, fill=black!8] (s) {\textbf{Spór między akcjonariuszami}};
  \node[box=3.0cm, right=of s] (n) {negocjacje\\30 dni od wezwania};
  \node[xx=3.0cm, right=of n] (m) {mediacja\\(ośrodek z UA)};
  \node[xx=3.0cm, right=of m] (sd) {sąd powszechny\\siedziby Spółki};
  \draw[line] (s) -- (n); \draw[line] (n) -- (m); \draw[line] (m) -- (sd);
  \node[wz=7.2cm, below=9mm of n.south east, anchor=north] (imp) {\textbf{Impas:} WZ dwa razy (w odstępie $\ge$30 dni) nie podejmuje uchwały o emisji (lit.~b) albo o transakcji ponad 500\,000 zł (lit.~h), a sprawa jest niezbędna --- każdy może wszcząć mediację};
  \draw[line] (imp.north) ++(1.7,0) -- (m.south);
\end{tikzpicture}
\caption{Rozwiązywanie sporów (UA §~11). Zabezpieczenie sądowe jest dostępne zawsze.}
\end{figure}

% ================================================================
\section{Kalendarz terminów}

\begin{longtable}{@{}>{\raggedright\arraybackslash}p{6.2cm} >{\raggedright\arraybackslash}p{4.4cm} >{\raggedright\arraybackslash}p{3.6cm}@{}}
\toprule
\textbf{Czynność} & \textbf{Termin} & \textbf{Źródło} \\
\midrule
\endhead
Uchwała akcjonariuszy poza WZ & 3 dni robocze -- 21 dni na głos & US §~12 ust.~3 \\
Zgoda albo odmowa zgody na zbycie akcji & 14 dni od zgłoszenia & US §~8 ust.~1 \\
Dodatkowy termin po wezwaniu Rady & 7 dni & US §~8 ust.~1 \\
Uchwała akcjonariuszy w sprawie zgody & 14 dni od rozesłania projektu & US §~8 ust.~1 \\
Wskazanie nabywcy po odmowie (konkurent) & 30 dni od odmowy & US §~8 ust.~3 \\
Lock-up Założycieli Pierwotnych / serii F & 12 mies. od wpisu / Dnia Przyznania & US §~8 ust.~4 \\
Stwierdzenie ustania zaangażowania & 30 dni od dowiedzenia się & UA §~4 ust.~2 \\
Wskazanie nabywcy akcji niezwolnionych & 3 mies. od ustania & US §~9 ust.~4 \\
Nabycie przez pozostałych Założycieli & kolejne 30 dni & US §~9 ust.~4 \\
Wygaśnięcie obowiązku zbycia & 6 mies. od ustania & US §~9 ust.~4 \\
Zapłata ceny akcji niezwolnionych & 14 dni od umowy zbycia & UA §~5 ust.~4 \\
Wydanie utworów po odejściu & 14 dni & Umowa nr~2 §~8, polityka §~7 \\
Odebranie dostępów & najpóźniej ostatni dzień; przy sporze od razu & polityka §~7 \\
Zgłoszenie incydentu / naruszenia danych & 24 godziny / 72 godziny do UODO & polityka §~6 \\
Waloryzacja kar i Cennika & co roku 1 marca & UA §~9 ust.~2, Cennik §~5 \\
Negocjacje przed mediacją & 30 dni od wezwania & UA §~11 ust.~1 \\
Kadencja Rady Dyrektorów & 3 lata obrotowe & US §~10 ust.~3 \\
Termin emisji serii F i P & do 31 grudnia 2036 r. & US §~5 ust.~3 \\
\bottomrule
\end{longtable}

\end{document}
